\begin{ExerciseBlock}
\ExerciseHeading{03}{Ausmessen des Amplitudengangs}{forced_oscillations_amplitude_response}
Links ist der Amplitudengang $A(\Omega)$ eines viskos gedämpften Feder-Masse-Schwingers dargestellt, rechts ein Ausschnitt der Resonanzkurve. Lesen Sie die benötigten Daten aus den Diagrammen ab.
\begin{center}
\CachedTikzFigure{content/resources/chapters/forced_oscillations/tikz/amplitude_response.tex}
\end{center}
\begin{enumerate}
\item Bestimmen Sie die Güte $Q$ auf zwei Arten.
\item Welcher Bruchteil der Anfangsamplitude bleibt nach Abschalten des Erregers nach fünf Perioden übrig? Wie lange dauern diese fünf Perioden?
\item Auf welchen Wert müsste man den Reibungskoeffizienten $b$ einstellen, damit bei $m=\qty{1.253}{\kilogram}$ der aperiodische Grenzfall eintritt?
\item Wie groß ist die Federkonstante?
\end{enumerate}
\end{ExerciseBlock}
\UseExerciseSolution{forced_oscillations_amplitude_response}
